\chapter{Localization of a product of fields at a single-coordinate element}
\label{chap:Mathlib_RingTheory_Localization_AtIdempotent}

This chapter tracks the Mathlib-PR-target helper
\lean{Localization.Away_pi_field_supportedAt}
in \texttt{Proetale/Mathlib/RingTheory/Localization/AtIdempotent.lean}.
It is the residue-identification ingredient needed to close the §5 Gap A
sorry of
\texttt{Proetale/Mathlib/RingTheory/Etale/StrictlyHenselian.lean}
(see chapter \texttt{local-structure.tex},
\texttt{lem:etale-locally-standard-etale-at-max}).

The result is folklore --- there is no external citation to quote
verbatim; this is project-bespoke infrastructure that fills a Mathlib
gap.

\begin{theorem}
  \label{thm:localization-away-pi-field-supported-at}
  \lean{Localization.Away_pi_field_supportedAt}
  \leanok
  Let $S$ be a commutative ring, $I$ a finite type with decidable
  equality, and $\{k_i\}_{i \in I}$ a family of fields each carrying
  an $S$-algebra structure. Fix $i_0 \in I$ and an element
  $r = (r_i)_{i \in I} \in \prod_{i \in I} k_i$ supported only at
  $i_0$: that is, $r_{i_0} \neq 0$ and $r_i = 0$ for $i \neq i_0$.
  Then the localization $(\prod_{i \in I} k_i)[1/r]$ is canonically
  isomorphic, as an $S$-algebra, to the single factor $k_{i_0}$.
\end{theorem}
\begin{proof}
  \uses{thm:localization-away-pi-field-supported-at}
  \leanok
  Write $R := \prod_{i \in I} k_i$ and let $\pi : R \to k_{i_0}$ be the
  projection $x \mapsto x_{i_0}$. Both $R \to k_{i_0}$ and $S \to k_{i_0}$
  give $\pi$ an $S$-algebra structure compatible with the product
  structure on $R$.

  Since $\pi(r) = r_{i_0} \neq 0$ in the field $k_{i_0}$, $\pi(r)$ is a
  unit. Hence powers of $r$ map to units under $\pi$, and the universal
  property of the localization $R[1/r]$ produces a unique $S$-algebra
  homomorphism $\widetilde{\pi} : R[1/r] \to k_{i_0}$ with
  $\widetilde{\pi}(x/1) = \pi(x)$ for $x \in R$.

  We claim $\widetilde{\pi}$ is bijective.

  \emph{Surjectivity.} $\pi : R \to k_{i_0}$ is already surjective
  (given $y \in k_{i_0}$, take $x$ with $x_{i_0} = y$ and $x_j = 0$ for
  $j \neq i_0$). The factor $1/r$ does not obstruct surjectivity.

  \emph{Injectivity.} Suppose $\widetilde{\pi}(x/r^n) = 0$. Then
  $\pi(x) = x_{i_0} = 0$ in $k_{i_0}$. We must show $x/r^n = 0$ in
  $R[1/r]$, i.e.\ that there is some $m \geq 0$ with $r^m x = 0$ in $R$.
  Take $m = 1$. For every $j \in I$:
  \begin{itemize}
    \item if $j = i_0$, $(r \cdot x)_j = r_{i_0} \cdot x_{i_0} =
      r_{i_0} \cdot 0 = 0$;
    \item if $j \neq i_0$, $(r \cdot x)_j = r_j \cdot x_j = 0 \cdot x_j
      = 0$.
  \end{itemize}
  Hence $r \cdot x = 0$ in $R$, so $x/r^n = (r \cdot x)/r^{n+1} = 0$
  in $R[1/r]$.

  Combining surjectivity and injectivity, $\widetilde{\pi}$ is an
  $S$-algebra isomorphism $R[1/r] \xrightarrow{\sim} k_{i_0}$.
\end{proof}

\medskip\noindent\textbf{Consumer.} The lemma is invoked inside
\texttt{surjective\_standardEtalePair\_lift\_of\_etale} (Helper B.3.b
of the §5 close of
\texttt{lem:etale-locally-standard-etale-at-max}) to identify the
residue field $B_b / \mathfrak{m} \cdot B_b$ with the $i_n$-factor
$k_{i_n}$ of the étale-product decomposition $k \otimes_A B \simeq
\prod_{i} k_i$, where the inverted element $b = e \cdot f'(\beta)
\cdot s_B$ has its residue supported at $i_n$ (the index corresponding
to the maximal ideal $\mathfrak{n}$, via the isolating idempotent
$e$ of
\texttt{exists\_idempotent\_lift\_isolating\_at\_maximal}).
